\SourceLeaf{008}
\begin{center}\small
EIN NEUER BEWEIS DES SATZES,\\[5pt]
\textbf{DASS JEDES POLYNOM}\\[5pt]
IN EINER VARIABLEN\\[5pt]
IN REELLE FAKTOREN ERSTEN ODER ZWEITEN GRADES ZERLEGT WERDEN KANN.\\[8pt]
\rule{1.5cm}{0.3pt}
\end{center}
\Article{1}
Jede bestimmte algebraische Gleichung lässt sich auf die Form
\[
x^m+Ax^{m-1}+Bx^{m-2}+\text{usw.}+M=0
\]
bringen, wobei $m$ eine positive ganze Zahl ist. Bezeichnen wir die linke Seite dieser Gleichung mit $X$ und nehmen wir an, dass die Gleichung $X=0$ durch mehrere ungleiche Werte von $x$ erfüllt wird, etwa indem wir $x=\alpha$, $x=\beta$, $x=\gamma$ usw. setzen, so ist die Funktion $X$ durch das Produkt aus den Faktoren $x-\alpha$, $x-\beta$, $x-\gamma$ usw. teilbar. Umgekehrt: Wenn das Produkt aus mehreren linearen Faktoren $x-\alpha$, $x-\beta$, $x-\gamma$ usw. die Funktion $X$ teilt, so wird die Gleichung $X=0$ erfüllt, wenn man $x$ irgendeiner der Größen $\alpha$, $\beta$, $\gamma$ usw. gleichsetzt. Ist schließlich $X$ gleich einem Produkt aus $m$ solchen linearen Faktoren (ob sie nun alle verschieden oder einige von ihnen identisch sind), so können keine anderen linearen Faktoren die Funktion $X$ teilen. Eine Gleichung vom Grad $m$ kann daher nicht mehr als $m$ Wurzeln haben. Zugleich ist aber klar, dass eine Gleichung vom Grad $m$ \textit{weniger} Wurzeln haben kann, auch wenn sich $X$ in $m$ lineare Faktoren zerlegen lässt: Sind nämlich einige dieser Faktoren identisch, so ist die Anzahl der verschiedenen Möglichkeiten, die Gleichung zu erfüllen, notwendig kleiner als $m$. Der Einheitlichkeit halber haben es die Mathematiker dennoch vorgezogen, auch in diesem Fall zu sagen, die Gleichung habe $m$ Wurzeln und nur einige davon seien einander gleich; was sie sich durchaus erlauben können.
